% Source: Wald GR, physical PDF p. 313
%         (printed p. 300).
% Coverage: Conformal diagram of the spherical-collapse spacetime.
% Figures/tables/footnotes: Figure 12.1; no tables or footnotes.
% Status: source-reviewed and compile-clean.
% Uncertainties: none.

\begingroup
\renewcommand{\thefigure}{12.1}
\begin{figure}[H]
  \centering
  \begin{tikzpicture}[x=.85cm,y=.85cm,>=Latex,line cap=round,line join=round,font=\small]
    % Cylindrical conformal compactification.
    \draw[thick] (-1.35,-2.55) -- (-1.35,2.85);
    \draw[thick] ( 1.35,-2.55) -- ( 1.35,2.85);
    \draw[thick] (-1.35,2.85) arc[start angle=180,end angle=360,x radius=1.35,y radius=.28];
    \draw[thick] (-1.35,2.85) arc[start angle=180,end angle=0,x radius=1.35,y radius=.28];
    \draw[thick] (-1.35,-2.55) arc[start angle=180,end angle=360,x radius=1.35,y radius=.28];
    \draw[dashed] (-1.35,-2.55) arc[start angle=180,end angle=0,x radius=1.35,y radius=.28];
    % The physical spacetime closes at i^0; the dashed curves are the hidden
    % halves of the same null hypersurfaces on the cylindrical diagram.
    \coordinate (bottom) at (0,-2.02);
    \coordinate (singL) at (-.28,1.78);
    \coordinate (singR) at (.28,1.78);
    \coordinate (izero) at (.88,.20);
    \draw[thick] (singR) .. controls (.92,1.42) and (1.25,.72) .. (izero);
    \draw[dashed] (singL) .. controls (-1.10,1.25) and (-1.08,.45) .. (izero);
    \draw[thick] (bottom) .. controls (.72,-1.52) and (1.23,-.62) .. (izero);
    \draw[dashed] (bottom) .. controls (-.92,-1.36) and (-.92,-.25) .. (izero);
    \node[left] at (-1.42,1.30) {\(\mathcal I^+\)};
    \node[right] at (1.35,1.18) {\(\mathcal I^+\)};
    \node[left] at (-1.40,-.98) {\(\mathcal I^-\)};
    \node[right] at (1.35,-.72) {\(\mathcal I^-\)};
    \fill (izero) circle (1.25pt) node[right] {\(i^0\)};
    % The event horizon and r=0 line meet at the regular center in the past and
    % bound the narrow black-hole region II beneath the spacelike singularity.
    \path[fill=black!10] (bottom) -- (singL) -- (singR) -- cycle;
    \draw[thick] (bottom) .. controls (-.18,-.42) and (-.26,.84) .. (singL);
    \draw[thick] (bottom) -- (singR);
    \draw[thick,decorate,decoration={zigzag,segment length=2.4pt,amplitude=.75pt}]
      (singL) -- (singR);
    \foreach \yy in {-1.55,-1.25,-.95,-.65,-.35,-.05,.25,.55,.85,1.15,1.45}
      \draw[thin] ({-.28*(\yy+2.02)/3.8},\yy) -- ({.28*(\yy+2.02)/3.8},\yy);
    \node[fill=white,inner sep=1pt] at (-.03,1.18) {II};
  \end{tikzpicture}
  \caption{A conformal diagram of the same spacetime as shown in Figures~6.11
  and~6.12. From this conformal diagram, it is apparent that region II of the physical
  spacetime lies outside of \(J^-(\mathcal I^+)\). In contrast, in Figure~11.1,
  \(J^-(\mathcal I^+)\) includes the entire physical spacetime.}
  \label{fig:12.1}
\end{figure}
\endgroup
